\section{Discussion}
\label{sec:discussion}

\subsection{Industrial implications}

Three results have direct operational consequences. First, the drivers of tube life are the firing intensity, the steam-to-carbon
ratio and the excess air, in that order, and they act almost independently (Section~\ref{sec:rq1}). Of these, excess air is the
least appreciated: at fixed fuel it lowers the hot-spot temperature by shifting heat release down the tube, and the optimiser
exploits it in both directions, choosing minimum excess air when hydrogen is the priority and maximum excess air when life is. Excess
air is also the cheapest variable to move on an operating plant, and its effect can be verified with the existing pyrometer
survey \todo{comment on the flue-gas oxygen and efficiency penalty of higher excess air}.

Second, the trade-off between production, fuel and life has a knee that is far from the current operating point of the
calibrated plant. The base is dominated by a substantial part of the front, and the knee regime obtains more hydrogen at less fuel
per kmol and a fraction of the life cost, mainly by running at higher load with lower specific firing and the highest inlet
temperature. Whether these gains are realisable depends on constraints outside the twin (preheat-coil metallurgy at \SI{900}{K}
inlet temperature, PSA capacity at \SI{110}{\percent} load, downstream steam balance), and the steam-credit sweep shows that the
recommended steam-to-carbon ratio moves from 4.0 to about 2.7 depending on how the plant values steam. The twin cannot decide this,
but it makes the dependence explicit, which is the useful output: the plant's steam balance, not the reformer, determines whether
running steam-rich is a life-saving measure or a fuel penalty.

Third, the life cost of load-following is small compared with the life cost of catalyst ageing, provided the outlet temperature is
controlled. The practical message is that a plant that must follow demand should not fear the tube-life consequences of smooth load
changes, but should watch the end-of-campaign hot-spot temperature under a methane-slip target, where the twin predicts the fourth
year to cost 2.3 times the first per kmol of hydrogen. Under uncertainty this ageing conclusion is the most robust of the study,
and the conclusion on load-following the least. The difference between the two comes from the size of the nominal effect
(a factor of 2.3 against \SI{2}{\percent}), not from a difference in modelling quality, and it is an argument for reporting
robustness fractions alongside nominal results in any digital-twin study \todo{soften or strengthen this recommendation}.

\subsection{The flame-length response to load as a measurement need}

The one structural assumption that changes a conclusion is how the heat-release profile of the burners responds to the firing rate.
The calibrated model uses a single dimensionless flame length, and the cross-validation residuals at the upper peep hole
(Section~\ref{sec:verification}) already show that this is the model's weakest element: the partially loaded Plant C2 is predicted
too hot at the top and the fully loaded Plant A too cold. When the flame is allowed to shorten with load, the flexible-operation
conclusion reverses because the hot spot at reduced load moves into a hotter part of the furnace. The three published operating
points cannot separate the two hypotheses. What would settle the question is a pyrometer survey at three or four elevations at two
or three firing rates on one furnace, which most plants already perform in part during turnaround planning; a digital twin of the
kind presented here turns such a survey into an estimate of the flame-length response with its uncertainty, and the robustness
fractions would then no longer straddle 0.5. We regard this as the single most valuable measurement for tube-life-aware operation of
top-fired reformers \todo{check against burner-manufacturer data on flame length versus firing rate for the burner type}.

\subsection{Limitations}

The model represents one average tube. Tube-to-tube variation in a top-fired furnace, caused by burner imbalance, row position and
flow maldistribution, is known to be of the order of \SI{20}{K} in wall temperature \todo{cite a survey of tube-to-tube wall
temperature spread}, which by the Larson--Miller relation corresponds to a factor of two to three in rupture time between the
hottest and the average tube. The measurement group of the uncertainty analysis ($\pm 10$ K) partly covers this, and a per-tube
extension would replace the single wall-temperature offset by a distribution over the tube population, but the present results are
statements about an average tube.

The life model covers creep only. Thermal fatigue during start-up and shutdown, creep--fatigue interaction, carburisation and the
relaxation of thermal stresses across the wall are not modelled, and the hoop stress is the pressure stress alone; the thermal
stress from the radial temperature gradient, which is of comparable magnitude in thick tubes \todo{quantify the thermal stress for
this wall thickness and flux}, is neglected. The scenario results should therefore be read as the creep-life cost of the scenarios,
not their total life cost, and the ranking of flexible-operation scenarios could change if the number of thermal cycles were
counted.

The rupture curve is a placeholder. The manufacturer's minimum curve for a proprietary alloy reproduced by Yeh \cite{yeh2021} was
used because no public curve with scatter information was available in the transcribed form the module requires, and the absolute
rupture times it yields at the low hoop stress of this tube are not credible. The uncertainty analysis compensates with a
factor-of-two heat-to-heat scatter and a $\pm 1$ range on the Larson--Miller constant, and the creep group is the largest
contributor to the variance of $\log_{10} t_r$; replacing the curve by a heat-normalised fit to the NIMS data sheets for the
25Cr--35Ni--Nb alloys \cite{nims1990,nims1991} is the first step for any absolute statement. The module is built for this
replacement and the ingestion has been tested on synthetic data.

The scenarios are quasi-steady: each hour is a steady state of the twin, and the thermal inertia of the tube, the catalyst and the
furnace refractory is neglected. For load ramps of 0.1 per hour and events of one hour or longer this is a reasonable approximation
\todo{compare the tube and refractory thermal time constants with the ramp rates}, but it excludes the transient overshoot of the
wall temperature during a fast firing step, which would increase the cost of the over-firing events. Finally, the calibration rests
on three operating points of one furnace from a single published source, so the parameter intervals are conditional on that
furnace and the transfer to another reformer requires recalibration on its own survey data; the pipeline is designed so that this
is a matter of replacing one CSV file.
